
\needspace{4\baselineskip}
\color{uniblue}{\subsection{Сборка сигналов (СС)}}
\color{black}

\begin{enumerate}[label=\arabic{section}.\arabic{subsection}.\arabic*, labelsep=4pt, leftmargin=0pt, itemindent=57pt]

%%%%%%%%%%%%%%%%%%%%%%%%%%%%%%%%%%%%%%%%%%%%%%%%%%%%%%%%%%%%%%%%%%%%%%%%%% СБОРКА СИГНАЛОВ %%%%%%%%%%%%%%%%%%%%%%%%%%%%%%%%%%%%%%%%%%%%%%%%%%%%%%%%%%%%%%%%%%%%%%%
\item
Функциональный блок «СС» используется в качестве дополнительной логики для предупредительной сигнализации, которая описана в подразделе \ref{ps:punkt}.
\item
Функция «СС»  принимает различные сигналы от других функций и формирует следующие обобщенные сигналы:
\begin{itemize}
	\item общий сигнал срабатывания реле газовых защит «СС: ГЗ сигн»;
	\item общий сигнал неисправности изоляции цепей газовых защит «СС: Низ.изол. ГЗ»;
	\item общий сигнал блокировки цепей газовых защит «СС: ГЗ заблокирована»;
	\item общий сигнал срабатывания датчиков технологических защит «СС: ТЗ сигн»;
	\item общий сигнал неисправности изоляции цепей технологических защит «СС: Низ.изол. ТЗ»;
	\item общий сигнал блокировки цепей технологических защит «СС: ТЗ заблокирована»;	
	\item общий сигнал срабатывания датчиков технологической сигнализации «СС: ТС сигн»;
	\item общий сигнал неисправности оперативного тока газовых и технологических защит «СС: ОТ сигн»;	
	\item общий сигнал неисправности оперативного тока дуговых защит и УРОВ стороны НН «СС: ОТ НН сигн»;
	\item общий сигнал срабатывания дуговых защит и УРОВ стороны НН «СС: Внеш. откл»;
	\item сигнал нарушения положения переключателей в выходных цепях  «СС: Вых.цепи разобраны»;
	\item сигнал нарушения положения испытательных блоков «СС:БИ выведены»;
	\item общий сигнал срабатывания внешних функций «СС:Общ. внеш. сигнал».	
\end{itemize}

\item
Выходные сигналы функционального блока могут быть использованы для дальнейшего конфигурирования (например, на выходные реле) и применяться в цепях участковых и индивидуальных табло, сигнальных ламп, сигнальных реле, звуковой сигнализации и пр.

\item
Логика работы функционального блока «СС» выполнена в соответствии с рисунком \ref{ss:img1}. В таблице \ref{ss:tbl1} приведены параметры, необходимые для настройки функционального блока «СС». 

\small
\begin{longtable}{|>{\centering\arraybackslash}m{5.3cm}|>{\centering\arraybackslash}m{3.3cm}|>{\centering\arraybackslash}m{4.2cm}|>{\centering\arraybackslash}m{1.8cm}|>{\centering\arraybackslash}m{1cm}|}
\caption{Параметры для настройки функции «СС»\hfill\vspace{-0.5\baselineskip}}\label{ss:tbl1}\\ 
\hline
\rowcolor{gray!30}
Параметр (Параметр на ИЧМ) & Условное обозначение на схеме & Значение/ Диапазон & Единица измерения & Шаг \\ 
\hline
\endfirsthead
%\caption{\emph{Продолжение\hfill\vspace{-0.5\baselineskip}}} \\ % Отступ обозначения таблицы от таблицы и выравнивание надписи слева
\caption*{\hspace{3pt}\emph{Продолжение таблицы \ref{ss:tbl1}\hfill\vspace{-0.5\baselineskip}}} \\ % сделано по ГОСТ 2.105 п.6.8.7
\hline
\rowcolor{gray!30}
Параметр (Параметр на ИЧМ) & Условное обозначение на схеме & Значение/ Диапазон & Единица измерения & Шаг \\ 
%\hline
\endhead
%\hline
\endfoot
%\hline
\endlastfoot
%==+t1*GAPC1|DZT2_SignAssembly
%===t1*GAPC1
\centering Контроль положения SA1 (Контр\_полож\_SA1) & \centering SGF1 & \centering 0 = Не предусмотрено\\1 = Предусмотрено & \centering -- & \centering \arraybackslash -- \\
\hline
\centering Контроль положения SA2 (Контр\_полож\_SA2) & \centering SGF2 & \centering 0 = Не предусмотрено\\1 = Предусмотрено & \centering -- & \centering \arraybackslash -- \\
\hline
\centering Контроль положения SA3 (Контр\_полож\_SA3) & \centering SGF3 & \centering 0 = Не предусмотрено\\1 = Предусмотрено & \centering -- & \centering \arraybackslash -- \\
\hline
\centering Контроль положения SA4 (Контр\_полож\_SA4) & \centering SGF4 & \centering 0 = Не предусмотрено\\1 = Предусмотрено & \centering -- & \centering \arraybackslash -- \\
\hline
\centering Контроль положения SG1 (Контр\_полож\_SG1) & \centering SGF5 & \centering 0 = Не предусмотрено\\1 = Предусмотрено & \centering -- & \centering \arraybackslash -- \\
\hline
\centering Контроль положения SG2 (Контр\_полож\_SG2) & \centering SGF6 & \centering 0 = Не предусмотрено\\1 = Предусмотрено & \centering -- & \centering \arraybackslash -- \\
\hline
\centering Контроль положения SG3 (Контр\_полож\_SG3) & \centering SGF7 & \centering 0 = Не предусмотрено\\1 = Предусмотрено & \centering -- & \centering \arraybackslash -- \\
\hline
\centering Контроль опер. тока цепей ГЗ (Контр\_ОТ\_ГЗ) & \centering SGF8 & \centering 0 = Не предусмотрено\\1 = Предусмотрено & \centering -- & \centering \arraybackslash -- \\
\hline
\centering Контроль опер. тока цепей ТЗ, ТС (Контр\_ОТ\_ТЗ\_ТС) & \centering SGF9 & \centering 0 = Не предусмотрено\\1 = Предусмотрено & \centering -- & \centering \arraybackslash -- \\
\hline
\centering Контроль опер. тока цепей ЗДЗ НН1 (Контр\_ОТ\_ЗДЗ1) & \centering SGF10 & \centering 0 = Не предусмотрено\\1 = Предусмотрено & \centering -- & \centering \arraybackslash -- \\
\hline
\centering Контроль опер. тока цепей ЗДЗ НН2 (Контр\_ОТ\_ЗДЗ2) & \centering SGF11 & \centering 0 = Не предусмотрено\\1 = Предусмотрено & \centering -- & \centering \arraybackslash -- \\
\hline
\centering Контроль опер. тока цепей УРОВ НН1 (Контр\_ОТ\_УРОВ1) & \centering SGF12 & \centering 0 = Не предусмотрено\\1 = Предусмотрено & \centering -- & \centering \arraybackslash -- \\
\hline
\centering Контроль опер. тока цепей УРОВ НН2 (Контр\_ОТ\_УРОВ2) & \centering SGF13 & \centering 0 = Не предусмотрено\\1 = Предусмотрено & \centering -- & \centering \arraybackslash -- \\
\hline
%===t1
\end{longtable}
\normalsize

\vspace{2mm}
\begin{figure}[!h]
\centering
\includegraphics[width=0.95\textwidth,height=0.95\textheight,keepaspectratio]{\fbpath/41. Сборка сигн. ДЗТ2/_latex/img/ss.pdf}
\captionof{figure}{Функциональная схема «СС»}\label{ss:img1}
\end{figure}

\end{enumerate}
\FloatBarrier % Все float'ы ДО этой строки будут размещены здесь.